\documentclass[12pt]{scrartcl}
%packages used:
\usepackage[utf8]{inputenc} % Kodierung
%\usepackage[ngerman]{babel} % Sprache
\usepackage[english]{babel}
\usepackage{lmodern}
%\renewcommand*\familydefault{\sfdefault}
\usepackage{amsfonts}
\usepackage{graphicx}




\begin{document}
\thispagestyle{empty}
%Head
\begin{center}
\Large{Georg-August-Universität Göttingen}\\
\Large{Fakultät für Mathematik und Informatik}\\
\Large{Institut für Numerische und Angewandte Mathematik}
\end{center}
%vertical space 3cm
\vspace*{1.5cm}
\begin{center}
\includegraphics[scale=0.5]{siegel.png}
\end{center}
%middle part
\vspace*{1.5cm}
\begin{center}
\textbf{Python package:}\\
~\\
\textbf{\texttt{\Huge{itreg}}}\\
~\\
Iterative solvers for ill-posed problems\\
-\\
Documentation of the differences between Matlab and Python\\
\end{center}
%vertical space 8cm
\vspace*{2.5cm}
%bottom part
\begin{flushleft}
\begin{tabular}{llll}
\textbf{Authors:} & & Jonas Lueg&\\
& & Mariko Müller & \\
& & Henrik Wiechers & \\
& & \\
\textbf{Date:} & & \today &\\
\end{tabular}
\end{flushleft}

%table of contents
\newpage
\thispagestyle{empty}
\tableofcontents

%first page
\newpage
\pagenumbering{arabic}
\section{Introduction}
The toolbox provides methods to solve ill-conditioned non-linear systems of equations
$$T(x) = y.$$
Here  $T:~X \rightarrow Y$ is a non-linear mapping between finite dimensional spaces 
$$X = \mathbb{K}^n$$
and
$$Y = \mathbb{K}^m$$
where  $\mathbb{K}=\mathbb{R}$ or  $\mathbb{K}=\mathbb{C}$. $X$ and  $Y$ are equipped with inner products described by Gram matrices  $G_X$ and  $G_Y$:
$$\langle x_1,x_2 \rangle_X = x_1^\ast G_X x_2$$
and
$$\langle y_1,y_2 \rangle_Y = y_1^\ast G_Y y_2.$$
\section{Structure: Python versus Matlab}
\subsection{Structure of the Operators}
In Matlab, the operator $T$ with all its attributes is given by a certain structure. This structure also contains the spaces $X$ and $Y$ with their respective inner products and other attributes. In Python this structure is divided in classes. We have the classes \texttt{Space}, \texttt{Operator} and \texttt{Solver}. Using objects of class \texttt{Space}, objects from class \texttt{Operator} are comparable to the structure of the operator in Matlab.
\subsection{Structure of the Solvers}
An object from class \texttt{Solver} in Python is equivalent to the respective solver in Matlab. As the name of the package implies, all algorithms are iterative procedures. Thus, every solver in Matlab contains at least one while loop. This fact is used in the Python code. We use the built-in Python functions \verb|__iter__| and \texttt{next} to replace the outer while loop. \texttt{next} contains all computations that the outer while loop contains in the respective Matlab algorithm. If an algorithm has additional while loops, they are implemented in Python as while loops, too.
\subsection{Usage of the Solvers}
To execute an algorithm in Matlab, one has to call the algorithm directly with a suiting input, immediately beginning to compute and giving an output. Since in Python a solver is a class, one has to define a variable, e.g.\ \texttt{mysolver}, as an object of class \texttt{Solver}, e.g.\ \texttt{Landweber}. One would write\\
\texttt{mysolver = Landweber(input)}.\\
Nothing is computed yet, as it is only a definition. The execution of the solver is then accomplished by using the attribute \texttt{run}. This function is the same for all algorithms defined so far and it only takes an object of class \texttt{Stoprule} as an argument.\\
In Matlab stoprules are given through parameters to the function. In Python there exists the class \texttt{StopRule}. Thus, one defines a stoprule \texttt{mystoprule} which is then given as an argument:\\
\texttt{x = mysolver.run(stoprule = mystoprule)}.\\
Then the solution for the input \texttt{input}, using the solver \texttt{mysolver}, is given by \texttt{x}.
\subsection{Inner Solvers}
Inner solvers are implemented as a class \verb|Inner_Solver|, analogously to the class \texttt{Solver}. Since all inner solvers as well as solvers have  at least one while loop, we use the same structure with \verb|__iter__| and \texttt{next}. Inner solvers are only used inside solvers and thus are in general not explicitly called by the user.
\subsection{Parameter Structure}
To give parameters like tolerances or maximal number of iterations to an algorithm in Matlab, one has to define all parameters inside a certain parameter structure. There are also pre-set parameters given in a file. In Python all those parameters have to be given as arguments to the function that constructs an object of class \texttt{Solver}, e.g. \texttt{Landweber}. These are mostly optional arguments, as there are pre-set values for each parameter.
\subsection{Utilities}
The \texttt{utilities} folder in the Matlab package corresponds to the \texttt{util} folder in the Python package. All utilities that are needed so far in Python are saved here as methods.
\subsection{Logging}
For diagnostic messages in Matlab, the algorithms print the messages directly in the console. In Python we want to use the standard \texttt{logging} framework. But so far this is only implemented in the \texttt{Landweber} algorithm.
\section{Details}
In this section, some details of differences between the Python and the Matlab package are discussed.
\subsection{WeightedOp}
In Matlab in the inner solver \texttt{SQP}, which is used by \verb|IRNM_KL|, a new operator $T$ is defined by weighting the \texttt{adjoint} and the \texttt{derivative} method of the old operator $F$. These definitions are written directly as commands in the Matlab code.\\
In Python copying this style is called monkey patching, since we would replace attributes of functions at runtime. To solve this issue, a new class \texttt{WeightedOp} of class \texttt{Operator} is created, that takes an operator and a weight as input and returns the desired weighted operator.
\subsection{Quadprog}
In Matlab in the solver \verb|IRGNM_L1_fid|, the built-in function \texttt{quadprog} is used to minimize the term for $x\in\mathbb{K}^n$:
$$\frac{1}{2}x^THx+f^Tx,$$
such that
$$L\leq x\leq R,$$
where $H\in\mathbb{K}^{n\times n}$ and $f,L,R\in\mathbb{K}^n$ are given parameters.\\
In Python this built-in function does not exist. Therefore we use a function from the SciPy package:\\
\texttt{scipy.optimize.minimize}.\\
It takes a function 
$$f:~\mathbb{K}^n\rightarrow\mathbb{R}$$
and boundaries $L,R\in\mathbb{K}^n$ as an input and minimizes
$$f(x),\quad x\in\mathbb{K}^n,~L\leq x\leq R.$$
\subsection{Plots}
In contrast to the Matlab package, no plots are included in the Python code. Thus, the results have to be plotted manually.

\end{document}